\begin{proof}
\begin{enumerate}[label={$\langle1\rangle$\arabic*.},leftmargin=*]
\item 
\noindent\textbf{Step 1.} The inclusions \(\mathcal O_Y(-(k+1)D)\hookrightarrow
\mathcal O_Y(-kD)\) give a descending chain
\[
 V_k=H^0(Y,B\otimes\mathcal O_Y(-kD))
 \subseteq H^0(Y,B).
\]
It stabilizes because the last vector space is finite-dimensional.  A
section in the stable value belongs to every \(V_k\).  In a local equation
\(z=0\) for \(D\), every coefficient of that section is divisible by
\(z^k\) for all \(k\), and hence is zero.  The section vanishes on a
neighbourhood of \(D\), and therefore on connected \(Y\) by analytic
continuation.  Thus \(V_k=0\) for some \(k\).  Since
\(\mathcal O_Y(-kD)\in\operatorname{Pic}^0(Y)\), upper semicontinuity in a
local universal family of topologically trivial line bundles shows that
\[
 \{P:h^0(Y,B\otimes P)>0\}
\]
is a proper analytic subset of \(\operatorname{Pic}^0(Y)\).  This proves
the first vanishing for general \(P\).  Serre duality identifies the dual
of the second group in \eqref{wbII:eq:W:5.23} with
\[
 H^0(Y,B^\vee\otimes\omega_Y\otimes P^\vee).
\]
Apply the first part to \(B^\vee\otimes\omega_Y\) and use that inversion is
an automorphism of \(\operatorname{Pic}^0(Y)\).

\end{enumerate}
\end{proof}
